\hypertarget{mux1ee5c-tiuxeau-tuux1ea7n-3}{%
\subsubsection{Mục tiêu tuần 3:}}

\begin{itemize}
\tightlist
\item
  Phát triển hoàn chỉnh \texttt{NewsRAGSpider} --- spider chính để crawl
  tin tức.
\item
  Xử lý logic extract thông tin bài viết: title, content, author,
  publish\_date.
\item
  Viết \texttt{KafkaPipeline} để đẩy dữ liệu crawl được vào Kafka topic.
\item
  Test crawl thử trên local.
\end{itemize}

\hypertarget{cuxe1c-cuxf4ng-viux1ec7c-cux1ea7n-triux1ec3n-khai-trong-tuux1ea7n-nuxe0y}{%
\subsubsection{Các công việc cần triển khai trong tuần
này:}}

\begingroup
\small
\setlength{\tabcolsep}{3pt}
\renewcommand{\arraystretch}{1.15}
\begin{longtable}[]{@{}
  >{\raggedright\arraybackslash}p{(\columnwidth - 8\tabcolsep) * \real{0.0119}}
  >{\raggedright\arraybackslash}p{(\columnwidth - 8\tabcolsep) * \real{0.7194}}
  >{\raggedright\arraybackslash}p{(\columnwidth - 8\tabcolsep) * \real{0.0474}}
  >{\raggedright\arraybackslash}p{(\columnwidth - 8\tabcolsep) * \real{0.0593}}
  >{\raggedright\arraybackslash}p{(\columnwidth - 8\tabcolsep) * \real{0.1621}}@{}}
\toprule\noalign{}
\begin{minipage}[b]{\linewidth}\raggedright
Thứ
\end{minipage} & \begin{minipage}[b]{\linewidth}\raggedright
Công việc
\end{minipage} & \begin{minipage}[b]{\linewidth}\raggedright
Ngày bắt đầu
\end{minipage} & \begin{minipage}[b]{\linewidth}\raggedright
Ngày hoàn thành
\end{minipage} & \begin{minipage}[b]{\linewidth}\raggedright
Nguồn tài liệu
\end{minipage} \\
\midrule\noalign{}
\endhead
\bottomrule\noalign{}
\endlastfoot
2 & - Viết class \texttt{NewsRAGSpider} kế thừa \texttt{scrapy.Spider} -
Implement \texttt{parse()} method: lọc URL nội bộ, phân loại link bài
viết (.html, .htm) và link danh mục - Cấu hình
\texttt{custom\_settings}: CONCURRENT\_REQUESTS, DOWNLOAD\_DELAY,
DEPTH\_LIMIT & 16/06/2026 & 16/06/2026 & spider.py \\
3 & - Implement \texttt{parse\_article()}: sử dụng \texttt{newspaper3k}
để parse HTML - Lọc bài viết có nội dung \textless{} 100 ký tự - Extract
author từ \texttt{article.authors} & 17/06/2026 & 17/06/2026 & \\
4 & - Xử lý author extraction nâng cao: \qquad  + Fallback qua nhiều CSS
selectors (\texttt{.author-name}, \texttt{.tac-gia},
\texttt{a{[}rel="author"{]}},\ldots) \qquad  + Validate author: loại bỏ
fake authors (URLs, ngày tháng, tên báo) \qquad  + Hàm
\texttt{is\_valid\_author()}: kiểm tra độ dài, ký tự đặc biệt, bad words
& 18/06/2026 & 18/06/2026 & \\
5 & - Xử lý publish\_date parsing: \qquad  + Parse ISO format:
\texttt{2026-07-01T14:30} \qquad  + Parse VN format:
\texttt{17/06/2026\ 14:30} \qquad  + Fallback qua nhiều CSS selectors +
\texttt{article.publish\_date} - Viết \texttt{KafkaPipeline}: serialize
item thành JSON \(\rightarrow\) push vào Kafka topic \texttt{news\_raw}
& 19/06/2026 & 19/06/2026 & pipelines.py \\
6 & - Test crawl local: \texttt{scrapy\ crawl\ news\_rag\_spider} - Kiểm
tra output: đúng format JSON, author hợp lệ, date đúng - Debug và sửa
lỗi CSS selectors cho từng báo & 20/06/2026 & 20/06/2026 & \\
\end{longtable}
\endgroup

\hypertarget{kux1ebft-quux1ea3-ux111ux1ea1t-ux111ux1b0ux1ee3c-tuux1ea7n-3}{%
\subsubsection{Kết quả đạt được tuần
3:}}

\begin{itemize}
\tightlist
\item
  Hoàn thành \texttt{NewsRAGSpider} (\texttt{crawler/spiders/spider.py})
  với các tính năng:

  \begin{itemize}
  \tightlist
  \item
    \textbf{parse()}: Duyệt tất cả link trên trang, lọc URL cùng domain,
    phân loại bài viết và trang danh mục
  \item
    \textbf{parse\_article()}: Sử dụng \texttt{newspaper3k} kết hợp CSS
    Selectors
  \item
    \textbf{Custom settings}: CONCURRENT\_REQUESTS=16 (Windows) / 32
    (Linux), DOWNLOAD\_DELAY=1.0/0.5, DEPTH\_LIMIT=5
  \end{itemize}
\item
  Xây dựng hệ thống extract author phức tạp với multi-level fallback:

  \begin{enumerate}
  \def\labelenumi{\arabic{enumi}.}
  \tightlist
  \item
    Lấy từ \texttt{newspaper3k} \(\rightarrow\) validate
  \item
    Fallback qua 12+ CSS selectors: \texttt{.author-name},
    \texttt{.tac-gia}, \texttt{a{[}href*="tac-gia"{]}},\ldots{}
  \item
    Scan bottom paragraphs của bài viết tìm pattern ``Theo: \ldots{}'' /
    ``Nguon: \ldots{}''
  \item
    Cuối cùng fallback về \texttt{meta{[}name="author"{]}}
  \end{enumerate}

  \begin{itemize}
  \tightlist
  \item
    Hàm \texttt{is\_valid\_author()} kiểm tra: độ dài 2-100 ký tự, không
    chứa URL/email, không chứa ngày tháng, không phải bad words (tên
    báo, nhãn chung chung)
  \end{itemize}
\item
  Xử lý date parsing đa dạng:

  \begin{itemize}
  \tightlist
  \item
    ISO format: \texttt{2026-07-01T14:30:00}
  \item
    VN format: \texttt{17/06/2026\ 14:30} hoặc \texttt{01-07-2026}
  \item
    Fallback qua 10 CSS selectors + \texttt{article.publish\_date}
  \item
    Output chuẩn: \texttt{YYYY-MM-DD\ HH:MM:SS}
  \end{itemize}
\item
  Hoàn thành \texttt{KafkaPipeline} (\texttt{crawler/pipelines.py}):

  \begin{itemize}
  \tightlist
  \item
    Serialize item thành JSON (ensure\_ascii=False cho tiếng Việt)
  \item
    Push vào Kafka topic \texttt{news\_raw}
  \item
    Producer flush sau mỗi item
  \end{itemize}
\item
  Test crawl thành công trên local, output format hợp lệ
\end{itemize}
